\begin{kbvderiv}
\vspace{4pt plus 3pt minus 1pt}
\begin{center}\begin{adjustbox}{max width=\textwidth}
{\renewcommand{\arraystretch}{1.22}\setlength{\tabcolsep}{4pt}
\begin{tabular}{l c l}
 & {\color{kbvgrey}$=$} & ನಾವು + ಎಲ್ಲರೂ \\
 & {\color{kbvgrey}$=$} & ನಾ + (ವ್ + ಉ) + ಎಲ್ಲರೂ {[}\ensuremath{\because} ವು = ವ್ + ಉ{]} \\
 & {\color{kbvgrey}$=$} & ನಾ + (ವ್ + ಎ) + ಎಲ್ಲರೂ.... {[}ಉ-ಕಾರ ಲೋಪ{]} \\
 & {\color{kbvgrey}$=$} & ನಾವೆಲ್ಲರೂ ..... {[}\ensuremath{\because} ವ್ + ಎ = ವೆ{]} \\
\end{tabular}}
\end{adjustbox}\end{center}
\vspace{4pt plus 3pt minus 1pt}
\end{kbvderiv}

\indent ನಾವು, ಎಲ್ಲರೂ ಎಂಬ ಪದಗಳಲ್ಲಿ ವು, ಎ, ಅಕ್ಷರಗಳು ಒಂದಕ್ಕೊಂದು ಸೇರಿ ಸಂಧಿಯಾದಾಗ ಉ-ಕಾರವು ಕಾಣದಾಗುವುದು. ಇದನ್ನು ಉ-ಕಾರ ಲೋಪವಾಯಿತು ಎಂಬುದಾಗಿ ಹೇಳುತ್ತೇವೆ. ಹೀಗೆ ಸಂಧಿಕಾರ್ಯವು ನಡೆಯುವಾಗ ಒಂದು ಅಕ್ಷರವು ಬಿಟ್ಟು ಹೋಗುವುದಕ್ಕೆ “ಲೋಪ ಸಂಧಿ” ಎಂದು ಹೆಸರು.\par

\begin{kbvnote}
\noindent{\kbvsecfont ವಿವರಣೆ :}\hspace{0.5em}ಮೇಲಿನ ಉದಾಹರಣೆಯಲ್ಲಿ “ಉ” ಮತ್ತು “ಎ” ಸ್ವರಾಕ್ಷರ ಗಳು, ಹೀಗೆ, ಎರಡು ಸ್ವರಾಕ್ಷರಗಳು ಸಂಧಿಸುವಾಗ ಪದದ ಅರ್ಥ ಕೆಡದ ರೀತಿಯಲ್ಲಿ ಮೊದಲನೆಯ ಸ್ವರಾಕ್ಷರವು ಲೋಪವಾಗುವುದು. ಯಾವ ಅಕ್ಷರ ಲೋಪವಾಗುವುದೋ ಆ ಅಕ್ಷರವನ್ನು ಸೇರಿಸಿ ಸಂಧಿಯನ್ನು ಹೆಸರಿಸುವುದು ವಾಡಿಕೆಯಾಗಿದೆ. ಇಲ್ಲಿ ಉ-ಕಾರ ಲೋಪವಾಗಿರುವುದರಿಂದ, ಇದು ಉ-ಕಾರ ಲೋಪ ಸಂಧಿ.
\end{kbvnote}

\vspace{3pt plus 2pt}\noindent\hspace*{1.2em}\kbvlabel{ಉದಾ}\par\vspace{1pt}

\begin{kbvderiv}
\vspace{4pt plus 3pt minus 1pt}
\begin{center}\begin{adjustbox}{max width=\textwidth}
{\renewcommand{\arraystretch}{1.22}\setlength{\tabcolsep}{4pt}
\begin{tabular}{l c l}
ಅಲ್ಲಿ + ಅಲ್ಲಿ & {\color{kbvgrey}$=$} & \kbvresult{ಅಲ್ಲಲ್ಲಿ,}{\color{kbvgrey} .... (ಇ-ಕಾರ ಲೋಪ ಸಂಧಿ)} \\
ನನ್ನ + ಅಕೆ & {\color{kbvgrey}$=$} & \kbvresult{ನನ್ನಾಕೆ,}{\color{kbvgrey} .... (ಅ-ಕಾರ ಲೋಪ ಸಂಧಿ)} \\
ನನಗೆ + ಅಲ್ಲಿ & {\color{kbvgrey}$=$} & \kbvresult{ನನಗೆಲ್ಲಿ,}{\color{kbvgrey} .... (ಎ-ಕಾರ ಲೋಪ ಸಂಧಿ)} \\
\end{tabular}}
\end{adjustbox}\end{center}
\vspace{4pt plus 3pt minus 1pt}
\end{kbvderiv}

\vspace{4pt plus 4pt minus 1pt}\noindent{\kbvsecfont\fontsize{11pt}{14.5pt}\selectfont\color{kbvaccentdark} (೨)\hspace{0.45em}ಆಗಮ ಸಂಧಿ:}\hspace{0.5em}“ಕಳ್ಳರು ಕುರಿಯನ್ನು ಕದ್ದರು” ಎಂಬ ವಾಕ್ಯದಲ್ಲಿ ಕುರಿಯನ್ನು ಎಂಬ\par

\begin{kbvderiv}
\vspace{4pt plus 3pt minus 1pt}
\begin{center}\begin{adjustbox}{max width=\textwidth}
{\renewcommand{\arraystretch}{1.22}\setlength{\tabcolsep}{4pt}
\begin{tabular}{l c l}
ಸಂಧಿಪದವನ್ನು ಬಿಡಿಸಿ ನೋಡೋಣ. ಕುರಿ + ಅನ್ನು & {\color{kbvgrey}$=$} & \kbvresult{ಕುರಿಯನ್ನು} \\
ಕುರಿ + ಅನ್ನು & {\color{kbvgrey}$=$} & ಕು + (ರ್ + ಇ) + (ಅ) + ನ್ನು \\
ಕುರಿಯನ್ನು & {\color{kbvgrey}$=$} & ಕು + (ರ್ + ಇ + ಯ್ + ಅ) + ನ್ನು \\
\end{tabular}}
\end{adjustbox}\end{center}
\vspace{4pt plus 3pt minus 1pt}
\end{kbvderiv}

\indent ಇಲ್ಲಿ ಸ್ವರದ ಮುಂದೆ ಸ್ವರ ಬಂದು ಸಂಧಿಕಾರ್ಯವಾಗುವುದನ್ನು ಕಾಣ\par

\indent ಬಹುದು. ಹಿಂದೆ ತಿಳಿಸಿರುವಂತೆ ಸಂಧಿಮಾಡಿದಲ್ಲಿ ಕುರಿ + ಅನ್ನು = ಕುರೆನ್ನು\kbvcontinues
